\section{Incremental recursive programs}\label{sec:nested}

In \secref{sec:streams}--\ref{sec:relational} 
we showed how to incrementalize a relational query by
compiling it into a circuit, lifting the circuit to compute on streams, and
applying the $\inc{\cdot}$ operator to the lifted circuit.  In \secref{sec:datalog} we showed
how to compile a recursive query into a circuit that employs incremental
computation internally to compute the fixed point.
Here we combine these results to construct a circuit that evaluates a \emph{recursive
query incrementally}.  The circuit receives a stream of updates to input
relations, and for every update recomputes the fixed point.  To do this
incrementally, it preserves the stream of changes to recursive relations
produced by the iterative fixed point computation, and adjusts this stream to
account for the modified inputs.  Thus, every element of the input stream yields
a stream of adjustments to the fixed point computation, using
\emph{nested streams}.

Nested streams, or streams of streams, $\stream{\stream{A}} = \N \rightarrow (\N
\rightarrow A)$, are well defined, since streams form an abelian group.
Equivalently, a nested stream is a value in $\N \times \N \to A$, i.e., a matrix
with an infinite number of rows, indexed by two-dimensional time $(t_0, t_1)$. 
where each row is a stream.  In
\secref{sec-nested-examples} we show a few example nested stream
computations.

Lifting a stream operator $S: \stream{A} \to \stream{B}$ yields an operator over
nested streams $\lift{S}: \stream{\stream{A}} \to \stream{\stream{B}}$, such
that $(\lift{S})(s) = S \circ s$, or, pointwise: $(\lift{S}(s))[t_0][t_1] =
S(s[t_0])[t_1], \forall t_0, t_1 \in \N$.  In particular, a scalar function $f:
A \rightarrow B$ can be lifted twice to produce an operator between streams of
streams: $\lift{\lift{f}}: \stream{\stream{A}} \rightarrow \stream{\stream{B}}$.

We define a partial order over timestamps: $(i_0, i_1)
\leq (t_0, t_1)$ iff $i_0 \leq t_0$ and $i_1 \leq t_1$.  We extend the
definition of strictness for operators over nested streams: a stream operator
$F: \stream{\stream{A}} \to \stream{\stream{B}}$ is strict if for any $s, s' \in
\stream{\stream{A}}$ and all times $t, i \in \N \times \N$ we have $\forall i <
t, s[i] = s'[i]$ implies $F(s)[t] = F(s')[t]$.
Proposition~\ref{prop-unique-fix} holds for this notion of strictness, i.e., the fixed point operator $\fix{\alpha}{F(\alpha)}$ is well defined for a strict operator $F$.

\begin{proposition}\label{prop-liftz}
The operator $\lift{\zm}: \stream{\stream{A}} \to \stream{\stream{A}}$ is strict.
\end{proposition}

The operator $\zm$ on nested streams delays ``rows'' of the matrix, 
while $\lift{\zm}$ delays ``columns''.  (See examples in 
\secref{sec-nested-examples}).

The $\I$ operator on $\stream{\stream{A}}$ operates on rows
of the matrix, treating each row as a single value.
Lifting a stream operator computing on $\stream{A}$, 
such as $\I: \stream{A} \to \stream{A}$, also produces an operator on nested streams, but
this time computing on the columns of the matrix
$\lift{\I}: \stream{\stream{A}} \to \stream{\stream{A}}.$  

\begin{proposition}[Lifting cycles]
\label{prop-lift-cycle}
For a binary, causal $T$ we have:
$\lift{(\lambda s. \fix{\alpha}{T(s,\zm(\alpha)}))} = \lambda s. \fix{\alpha}{(\lift{T})(s,(\lift{\zm})(\alpha))}$
\noindent i.e., lifting a circuit containing a ``cycle'' can be accomplished by
lifting all operators independently, including the $\zm$ back-edge.
\end{proposition}

This means that lifting a \dbsp stream function can be expressed within \dbsp
itself.  For example, we have:

\begin{tabular}{m{2cm}m{.5cm}m{4cm}}
\begin{tikzpicture}[>=latex]
  \node[] (input) {$i$};
  \node[block, right of=input] (I) {$\lift{\I}$};
  \node[right of=I] (output)  {$o$};
  \draw[->] (input) -- (I);
  \draw[->] (I) -- (output);
\end{tikzpicture}
& $\cong$ &
\begin{tikzpicture}[>=latex]
  \node[] (input) {$i$};
  \node[block, circle, right of=input, inner sep=0cm] (p) {$+$};
  \node[right of=p, node distance=1.5cm] (output)  {$o$};
  \node[block, below of=p, node distance=.7cm] (z) {$\lift{\zm}$};
  \draw[->] (input) -- (p);
  \draw[->] (p) -- node (mid) {} (output);
  \draw[->] (z) -- (p);
  \draw[->] (mid.center) |- (z);
\end{tikzpicture}
\end{tabular}

This proposition gives the ability to lift
entire circuits, including circuits computing on streams and having feedback edges,
which are well-defined, due to Proposition~\ref{prop-liftz}.  
With this machinery we can now apply Algorithm~\ref{algorithm-inc} to arbitrary
circuits, even circuits built for recursively-defined relations.  
Consider the ``semi-naive'' circuit~(\ref{eq:seminaive}),
and denote $\distinct \circ R$ with $T$:

\begin{center}
\begin{tikzpicture}[>=latex]
  \node[] (Iinput) {\code{I}};
  \node[block, right of=Iinput] (Idelta) {$\delta_0$};
  \node[block, right of=Idelta] (f) {$\inc{(\lift{T})}$};
  \node[block, right of=f, node distance=1.5cm] (S) {$\int$};
  \node[right of=S] (output)  {\code{O}};
  \draw[->] (f) -- node (o) {} (S);
  \node[block, below of=o, node distance=.7cm] (z) {$\zm$};
  \draw[->] (Iinput) -- (Idelta);
  \draw[->] (S) -- (output);
  \draw[->] (o.center) -- (z);
  \draw[->] (z) -| (f);
  \draw[->] (Idelta) -- (f);
\end{tikzpicture}
\vspace{-2mm}
\end{center}

\noindent Lift the entire circuit using Proposition~\ref{prop-lift-cycle} and incrementalize it:

\begin{tikzpicture}[>=latex]
  \node[] (Iinput) {\code{I}};
  \node[block, right of=Iinput] (I) {$\I$};
  \node[block, right of=I] (Idelta) {$\lift{\delta_0}$};
  \node[block, right of=Idelta, node distance=1.5cm] (f) {$\lift{\inc{(\lift{T})}}$};
  \node[block, right of=f, node distance=1.5cm] (S) {$\lift{\int}$};
  \node[block, right of=S] (D) {$\D$};
  \node[right of=D] (output)  {\code{O}};
  \draw[->] (f) -- node (o) {} (S);
  \node[block, below of=o, node distance=.7cm] (z) {$\lift{\zm}$};
  \draw[->] (Iinput) -- (I);
  \draw[->] (I) -- (Idelta);
  \draw[->] (S) -- (D);
  \draw[->] (D) -- (output);
  \draw[->] (o.center) -- (z);
  \draw[->] (z) -| (f);
  \draw[->] (Idelta) -- (f);
\end{tikzpicture}

\noindent Now apply the chain rule to this circuit:
\begin{equation}
\vspace{-2.1ex}
\begin{aligned}
\label{eq:increcursive}
\begin{tikzpicture}[>=latex]
  \node[] (Iinput) {\code{I}};
  \node[block, right of=Iinput] (Idelta) {$\lift{\delta_0}$};
  \node[block, right of=Idelta, node distance=2cm] (f) {$\inc{(\lift{\inc{(\lift{T})}})}$};
  \node[block, right of=f, node distance=2cm] (S) {$\lift{\int}$};
  \node[right of=S] (output)  {\code{O}};
  \draw[->] (f) -- node (o) {} (S);
  \node[block, below of=o, node distance=.7cm] (z) {$\lift{\zm}$};
  \draw[->] (Iinput) -- (Idelta);
  \draw[->] (S) -- (output);
  \draw[->] (o.center) -- (z);
  \draw[->] (z) -| (f);
  \draw[->] (Idelta) -- (f);
\end{tikzpicture}
\end{aligned}
\end{equation}
This is the incremental version of an arbitrary recursive query.

\subsection{Example}\label{sec:recursive-example}

In this section we derive the incremental version of a circuit containing
recursion, by applying Algorithm~\ref{algorithm-inc}.  We start with a very simple
program, expressed in Datalog, which computes the transitive closure of a directed
graph:

\begin{lstlisting}[language=ddlog]
// Edge relation with head and tail
input relation E(h: Node, t: Node)
// Reach relation with source s and sink t
output relation R(s: Node, t: Node)
R(x, x) :- E(x, _).
R(x, x) :- E(_, x).
R(x, y) :- E(x, y).
R(x, y) :- E(x, z), R(z, y).
\end{lstlisting}

We haven't explained how Datalog is translated to circuits, but most Datalog
operators are relational in nature.  Assuming one could write 
recursive queries in SQL where a view is defined in terms of itself, the
above program would be implemented by the following (illegal) SQL query:

\begin{lstlisting}[language=SQL]
CREATE VIEW R AS 
(SELECT E.h, E.h FROM E)
UNION
(SELECT E.t, E.t FROM E)
UNION
(SELECT * FROM E)
UNION
(SELECT E.h, R.t 
 FROM E JOIN R 
 ON E.t = R.s)
\end{lstlisting}

We apply the algorithm from \refsec{sec:datalog} to create first the non-recursive circuit,
by assuming that \code{R} is already computed as a view \code{R1}, and using \code{R1}
in the definition of \code{R} instead of itself:

\begin{lstlisting}[language=SQL]
CREATE VIEW Reach AS 
(SELECT E.h, E.h FROM E)
UNION
(SELECT E.t, E.t FROM E)
UNION
(SELECT * FROM E)
UNION
(SELECT E.h, R1.t 
 FROM E JOIN R ON E.t = R1.s)
\end{lstlisting}

Now we implement this query as a \dbsp circuit with two inputs \code{E} and \code{R1}:

\noindent
\begin{tikzpicture}[>=latex, node distance=1.2cm]
  \node[] (E) {\code{E}};
  \node[above of=E, node distance=.6cm] (R1) {\code{R1}};
  \node[block, right of=R1] (j) {$\bowtie_{t=s}$};
  \node[block, right of=j] (pj) {$\pi_{h, t}$};
  \node[block, below of=j] (p1) {$\pi_{h}$};
  \node[block, right of=p1] (s1) {$\sigma_{h, h}$};
  \node[block, below of=p1, node distance=.6cm] (p2)  {$\pi_{t}$};
  \node[block, right of=p2] (s2) {$\sigma_{t, t}$};
  \node[below of=pj, node distance=.6cm] (mid) {};
  \node[block, circle, right of=mid, inner sep=0cm, node distance=1.5cm] (plus) {$+$};
  \node[block, right of=plus] (d) {$\distinct$};
  \node[right of=d] (R) {\code{R}};
  \draw[->] (R1) -- (j);
  \draw[->] (E) -- (j);
  \draw[->] (j) -- (pj);
  \draw[->] (E) -- (p1);
  \draw[->] (p1) -- (s1);
  \draw[->] (E) -- (p2);
  \draw[->] (p2) -- (s2);
  \draw[->] (E) -- (plus);
  \draw[->] (pj) -- (plus);
  \draw[->] (s1) -- (plus);
  \draw[->] (s2) -- (plus);
  \draw[->] (plus) -- (d);
  \draw[->] (d) -- (R);
\end{tikzpicture}

Now lift the circuit by lifting each operator pointwise, and connect it in a feedback
loop by connecting input \code{R1} from the output \code{R} through a \zm operator and bracket
everything with $\delta_0 -- \int$:

\noindent
\begin{tikzpicture}[>=latex]
  \node[] (Einput) {\code{E}};
  % generic part
  \node[block, right of=Einput, node distance=.8cm] (ID) {$\delta_0$};
  \node[block, right of=ID, node distance=.8cm] (E) {$\I$};
  
  % relational query
  \node[right of=E] (empty) {};
  \node[block, above of=empty, node distance=.6cm] (j) {$\lift{\bowtie_{t=s}}$};
  \node[block, right of=j, node distance=1.2cm] (pj) {$\lift{\pi_{h, t}}$};
  \node[block, below of=j] (p1) {$\lift{\pi_{h}}$};
  \node[block, right of=p1] (s1) {$\lift{\sigma_{h, h}}$};
  \node[block, below of=p1, node distance=.6cm] (p2)  {$\lift{\pi_{t}}$};
  \node[block, right of=p2] (s2) {$\lift{\sigma_{t, t}}$};
  \node[below of=pj, node distance=.6cm] (mid) {};
  \node[block, circle, right of=mid, inner sep=0cm, node distance=1.2cm] (plus) {$+$};
  \node[block, right of=plus] (d) {$\lift{\distinct}$};
  \draw[->] (E) -- (j);
  \draw[->] (j) -- (pj);
  \draw[->] (E) -- (p1);
  \draw[->] (p1) -- (s1);
  \draw[->] (E) -- (p2);
  \draw[->] (p2) -- (s2);
  \draw[->] (E) -- (plus);
  \draw[->] (pj) -- (plus);
  \draw[->] (s1) -- (plus);
  \draw[->] (s2) -- (plus);
  \draw[->] (plus) -- (d);
  
  % generic part
  \node[block, right of=d, node distance=1.1cm] (D) {$\D$};
  \node[block, right of=D, node distance=.8cm] (S) {$\int$};
  \node[right of=S, node distance=.8cm] (output)  {\code{R}};
  \draw[->] (Einput) -- (ID);
  \draw[->] (ID) -- (E);
  \draw[->] (d) -- node (o) {} (D);
  \draw[->] (D) -- (S);
  \draw[->] (S) -- (output);
  \node[block, above of=j, node distance=.8cm] (z) {$\zm$};
  \draw[->] (o.center) |- (z);
  \draw[->] (z) -- (j);
\end{tikzpicture}

The above circuit is a complete implementation of the non-streaming
recursive query; given an input relation \code{E} it will produce
its transitive closure \code{R} at the output.  

Now we use the semina\"ive property~\ref{eq:seminaive} to rewrite the circuit:

(To save space in the figures we will omit the indices from $\pi$ and $\sigma$
in the subsequent figures, for example by writing just $\pi$ instead of $\pi_h$.)

\noindent
\begin{tikzpicture}[>=latex, node distance=1.4cm]
  \node[] (Einput) {\code{E}};
  % generic part
  \node[block, right of=Einput, node distance=.8cm] (E) {$\delta_0$};
  
  % relational query
  \node[right of=E] (empty) {};
  \node[block, above of=empty, node distance=.6cm] (j) {$\inc{(\lift{\bowtie})}$};
  \node[block, right of=j] (pj) {$\inc{(\lift{\pi})}$};
  \node[block, below of=j, node distance=1cm] (p1) {$\inc{(\lift{\pi})}$};
  \node[block, right of=p1] (s1) {$\inc{(\lift{\sigma)}}$};
  \node[block, below of=p1, node distance=.6cm] (p2)  {$\inc{(\lift{\pi})}$};
  \node[block, right of=p2] (s2) {$\inc{(\lift{\sigma})}$};
  \node[below of=pj, node distance=.6cm] (mid) {};
  \node[block, circle, right of=mid, inner sep=0cm, node distance=1.2cm] (plus) {$+$};
  \node[block, right of=plus] (d) {$\inc{(\lift{\distinct})}$};
  \draw[->] (E) -- (j);
  \draw[->] (j) -- (pj);
  \draw[->] (E) -- (p1);
  \draw[->] (p1) -- (s1);
  \draw[->] (E) -- (p2);
  \draw[->] (p2) -- (s2);
  \draw[->] (E) -- (plus);
  \draw[->] (pj) -- (plus);
  \draw[->] (s1) -- (plus);
  \draw[->] (s2) -- (plus);
  \draw[->] (plus) -- (d);
  
  % generic part
  \node[block, right of=d] (S) {$\int$};
  \node[right of=S, node distance=.8cm] (output)  {\code{R}};
  \draw[->] (Einput) -- (E);
  \draw[->] (d) -- node (o) {} (S);
  \draw[->] (S) -- (output);
  \node[block, above of=j, node distance=.8cm] (z) {$\zm$};
  \draw[->] (o.center) |- (z);
  \draw[->] (z) -- (j);
\end{tikzpicture}

Using the linearity of $\lift\pi$ and $\lift\sigma$, this can be rewritten as an equivalent
circuit:

\noindent
\begin{tikzpicture}[>=latex, node distance=1.3cm]
  \node[] (Einput) {\code{E}};
  % generic part
  \node[block, right of=Einput, node distance=.8cm] (E) {$\delta_0$};
  
  % relational query
  \node[right of=E] (empty) {};
  \node[block, above of=empty, node distance=.6cm] (j) {$\inc{(\lift{\bowtie})}$};
  \node[block, right of=j] (pj) {$\lift{\pi}$};
  \node[block, below of=j, node distance=1cm] (p1) {$\lift{\pi}$};
  \node[block, right of=p1] (s1) {$\lift{\sigma}$};
  \node[block, below of=p1, node distance=.6cm] (p2)  {$\lift{\pi}$};
  \node[block, right of=p2] (s2) {$\lift{\sigma}$};
  \node[below of=pj, node distance=.6cm] (mid) {};
  \node[block, circle, right of=mid, inner sep=0cm, node distance=1.2cm] (plus) {$+$};
  \node[block, right of=plus] (d) {$\inc{(\lift{\distinct})}$};
  \draw[->] (E) -- (j);
  \draw[->] (j) -- (pj);
  \draw[->] (E) -- (p1);
  \draw[->] (p1) -- (s1);
  \draw[->] (E) -- (p2);
  \draw[->] (p2) -- (s2);
  \draw[->] (E) -- (plus);
  \draw[->] (pj) -- (plus);
  \draw[->] (s1) -- (plus);
  \draw[->] (s2) -- (plus);
  \draw[->] (plus) -- (d);
  
  % generic part
  \node[block, right of=d] (S) {$\int$};
  \node[right of=S, node distance=.8cm] (output)  {\code{R}};
  \draw[->] (Einput) -- (E);
  \draw[->] (d) -- node (o) {} (S);
  \draw[->] (S) -- (output);
  \node[block, above of=j, node distance=.8cm] (z) {$\zm$};
  \draw[->] (o.center) |- (z);
  \draw[->] (z) -- (j);
\end{tikzpicture}

To make this circuit into a streaming computation that evaluates a new transitive
closure for a stream of inputs \code{E}, we lift it entirely,
using Proposition~\ref{prop-lift-cycle}:

\noindent
\begin{tikzpicture}[>=latex, node distance=1.3cm]
  \node[] (Einput) {\code{E}};
  % generic part
  \node[block, right of=Einput, node distance=.8cm] (E) {$\lift{\delta_0}$};
  
  % relational query
  \node[right of=E] (empty) {};
  \node[block, above of=empty, node distance=.6cm] (j) {$\lift{\inc{(\lift{\bowtie})}}$};
  \node[block, right of=j] (pj) {$\lift{\lift{\pi}}$};
  \node[block, below of=j, node distance=1cm] (p1) {$\lift{\lift{\pi}}$};
  \node[block, right of=p1] (s1) {$\lift{\lift{\sigma}}$};
  \node[block, below of=p1, node distance=.6cm] (p2)  {$\lift{\lift{\pi}}$};
  \node[block, right of=p2] (s2) {$\lift{\lift{\sigma}}$};
  \node[below of=pj, node distance=.6cm] (mid) {};
  \node[block, circle, right of=mid, inner sep=0cm, node distance=1.2cm] (plus) {$+$};
  \node[block, right of=plus] (d) {$\lift{\inc{(\lift{\distinct})}}$};
  \draw[->] (E) -- (j);
  \draw[->] (j) -- (pj);
  \draw[->] (E) -- (p1);
  \draw[->] (p1) -- (s1);
  \draw[->] (E) -- (p2);
  \draw[->] (p2) -- (s2);
  \draw[->] (E) -- (plus);
  \draw[->] (pj) -- (plus);
  \draw[->] (s1) -- (plus);
  \draw[->] (s2) -- (plus);
  \draw[->] (plus) -- (d);
  
  % generic part
  \node[block, right of=d, node distance=1.5cm] (S) {$\lift{\int}$};
  \node[right of=S, node distance=.8cm] (output)  {\code{R}};
  \draw[->] (Einput) -- (E);
  \draw[->] (d) -- node (o) {} (S);
  \draw[->] (S) -- (output);
  \node[block, above of=j, node distance=.8cm] (z) {$\lift{\zm}$};
  \draw[->] (o.center) |- (z);
  \draw[->] (z) -- (j);
\end{tikzpicture}

We convert this circuit into an incremental circuit, which receives
in each transaction the changes to relation \code{E} and produces the 
corresponding changes to relation \code{R}:

\noindent
\begin{tikzpicture}[>=latex, node distance=1.3cm]
  \node[] (DE) {$\Delta$\code{E}};
  \node[block, right of=DE, node distance=.7cm] (Einput) {$\I$};
  \draw[->] (DE) -- (Einput);
  % generic part
  \node[block, right of=Einput, node distance=.8cm] (E) {$\lift{\delta_0}$};
  
  % relational query
  \node[right of=E, node distance=1.2cm] (empty) {};
  \node[block, above of=empty, node distance=.6cm] (j) {$\lift{\inc{(\lift{\bowtie})}}$};
  \node[block, right of=j] (pj) {$\lift{\lift{\pi}}$};
  \node[block, below of=j, node distance=1cm] (p1) {$\lift{\lift{\pi}}$};
  \node[block, right of=p1] (s1) {$\lift{\lift{\sigma}}$};
  \node[block, below of=p1, node distance=.6cm] (p2)  {$\lift{\lift{\pi}}$};
  \node[block, right of=p2] (s2) {$\lift{\lift{\sigma}}$};
  \node[below of=pj, node distance=.6cm] (mid) {};
  \node[block, circle, right of=mid, inner sep=0cm, node distance=1cm] (plus) {$+$};
  \node[block, right of=plus, node distance=1.2cm] (d) {$\lift{\inc{(\lift{\distinct})}}$};
  \draw[->] (E) -- (j);
  \draw[->] (j) -- (pj);
  \draw[->] (E) -- (p1);
  \draw[->] (p1) -- (s1);
  \draw[->] (E) -- (p2);
  \draw[->] (p2) -- (s2);
  \draw[->] (E) -- (plus);
  \draw[->] (pj) -- (plus);
  \draw[->] (s1) -- (plus);
  \draw[->] (s2) -- (plus);
  \draw[->] (plus) -- (d);
  
  % generic part
  \node[block, right of=d] (S) {$\lift{\int}$};
  \node[block, right of=S, node distance=.7cm] (OD) {$\D$};
  \node[right of=OD, node distance=.8cm] (output)  {$\Delta$\code{R}};
  \draw[->] (Einput) -- (E);
  \draw[->] (d) -- node (o) {} (S);
  \draw[->] (S) -- (OD);
  \draw[->] (OD) -- (output);
  \node[block, above of=j, node distance=.8cm] (z) {$\lift{\zm}$};
  \draw[->] (o.center) |- (z);
  \draw[->] (z) -- (j);
\end{tikzpicture}

We can now apply again the chain rule to this circuit:

\noindent
\begin{tikzpicture}[>=latex, node distance=1.4cm]
  \node[] (Einput) {$\Delta$\code{E}};
  % generic part
  \node[block, right of=Einput, node distance=1cm] (E) {$\inc{(\lift{\delta_0})}$};
  
  % relational query
  \node[right of=E] (empty) {};
  \node[block, above of=empty, node distance=.6cm] (j) {$\inc{(\lift{\inc{(\lift{\bowtie})}})}$};
  \node[block, right of=j, node distance=1.6cm] (pj) {$\inc{(\lift{\lift{\pi}})}$};
  \node[block, below of=j, node distance=1cm] (p1) {$\inc{(\lift{\lift{\pi}})}$};
  \node[block, right of=p1] (s1) {$\inc{(\lift{\lift{\sigma}})}$};
  \node[block, below of=p1, node distance=.6cm] (p2)  {$\inc{(\lift{\lift{\pi}})}$};
  \node[block, right of=p2] (s2) {$\inc{(\lift{\lift{\sigma}})}$};
  \node[below of=pj, node distance=.6cm] (mid) {};
  \node[block, circle, right of=mid, inner sep=0cm, node distance=1cm] (plus) {$+$};
  \node[block, right of=plus] (d) {$\inc{(\lift{\inc{(\lift{\distinct})}})}$};
  \draw[->] (E) -- (j);
  \draw[->] (j) -- (pj);
  \draw[->] (E) -- (p1);
  \draw[->] (p1) -- (s1);
  \draw[->] (E) -- (p2);
  \draw[->] (p2) -- (s2);
  \draw[->] (E) -- (plus);
  \draw[->] (pj) -- (plus);
  \draw[->] (s1) -- (plus);
  \draw[->] (s2) -- (plus);
  \draw[->] (plus) -- (d);
  
  % generic part
  \node[block, right of=d, node distance=1.8cm] (S) {$\inc{(\lift{\int})}$};
  \node[right of=S, node distance=1cm] (output)  {$\Delta$\code{R}};
  \draw[->] (Einput) -- (E);
  \draw[->] (d) -- node (o) {} (S);
  \draw[->] (S) -- (output);
  \node[block, above of=j, node distance=.8cm] (z) {$\inc{(\lift{\zm})}$};
  \draw[->] (o.center) |- (z);
  \draw[->] (z) -- (j);
\end{tikzpicture}

We now take advantage of the linearity of $\lift\delta_0$, $\lift\int$, 
$\lift\zm$, $\lift\lift\pi$, and $\lift\lift\sigma$ to simplify the circuit
by removing some $\inc{\cdot}$ invocations:

\noindent
\begin{tikzpicture}[>=latex, node distance=1.3cm]
  \node[] (Einput) {$\Delta$\code{E}};
  % generic part
  \node[block, right of=Einput, node distance=.8cm] (E) {$\lift{\delta_0}$};
  
  % relational query
  \node[right of=E] (empty) {};
  \node[block, above of=empty, node distance=.6cm] (j) {$\inc{(\lift{\inc{(\lift{\bowtie})}})}$};
  \node[block, right of=j, node distance=1.6cm] (pj) {$\lift{\lift{\pi}}$};
  \node[block, below of=j, node distance=1cm] (p1) {$\lift{\lift{\pi}}$};
  \node[block, right of=p1] (s1) {$\lift{\lift{\sigma}}$};
  \node[block, below of=p1, node distance=.6cm] (p2)  {$\lift{\lift{\pi}}$};
  \node[block, right of=p2] (s2) {$\lift{\lift{\sigma}}$};
  \node[below of=pj, node distance=.6cm] (mid) {};
  \node[block, circle, right of=mid, inner sep=0cm, node distance=.8cm] (plus) {$+$};
  \node[block, right of=plus] (d) {$\inc{(\lift{\inc{(\lift{\distinct})}})}$};
  \draw[->] (E) -- (j);
  \draw[->] (j) -- (pj);
  \draw[->] (E) -- (p1);
  \draw[->] (p1) -- (s1);
  \draw[->] (E) -- (p2);
  \draw[->] (p2) -- (s2);
  \draw[->] (E) -- (plus);
  \draw[->] (pj) -- (plus);
  \draw[->] (s1) -- (plus);
  \draw[->] (s2) -- (plus);
  \draw[->] (plus) -- (d);
  
  % generic part
  \node[block, right of=d, node distance=1.8cm] (S) {$\lift{\int}$};
  \node[right of=S, node distance=.8cm] (output)  {$\Delta$\code{R}};
  \draw[->] (Einput) -- (E);
  \draw[->] (d) -- node (o) {} (S);
  \draw[->] (S) -- (output);
  \node[block, above of=j, node distance=.8cm] (z) {$\lift{\zm}$};
  \draw[->] (o.center) |- (z);
  \draw[->] (z) -- (j);
\end{tikzpicture}

There are two applications of $\inc{\cdot}$ left in this circuit: $\inc{(\lift{\inc{(\lift{\bowtie})}})}$
and $\inc{(\lift{\inc{(\lift\distinct)}})}$.  We expand their implementations separately,
and we stitch them into the global circuit at the end.  This ability to reason about
sub-circuits independently highlights the modularity of \dbsp.

The join is expanded twice, using the bilinearity
of $\lift\bowtie$ and $\lift\lift\bowtie$.  Let's start with the inner circuit,
implementing $\inc{(\lift{\bowtie})}$, given by Theorem~\ref{bilinear}:

\begin{tabular}{m{2cm}m{.5cm}m{4.5cm}}
\begin{tikzpicture}[auto,>=latex]
    \node[] (a) {$a$};
    \node[below of=a, node distance=.6cm] (midway) {};
    \node[below of=midway, node distance=.6cm] (b) {$b$};
    \node[block, right of=midway] (q) {$\inc{(\lift{\bowtie})}$};
    \node[right of=q] (output) {$o$};
    \draw[->] (a) -| (q);
    \draw[->] (b) -| (q);
    \draw[->] (q) -- (output);
\end{tikzpicture} &
$\cong$ &
\begin{tikzpicture}[auto,>=latex]
  \node[] (a) {$a$}; 
  \node[block, below of=a, node distance=.6cm] (ab) {$\lift\bowtie$};
  \node[below of=ab, node distance=.6cm] (b) {$b$};
  \node[block, right of=a, node distance=1cm] (jI1) {$\I$};
  \node[block, right of=b, node distance=1cm] (jI2) {$\I$};
  \draw[->] (a) -- (jI1);
  \draw[->] (b) -- (jI2);
  \node[block, right of=jI1, node distance=.8cm] (ZI1) {$\zm$};
  \node[block, right of=jI2, node distance=.8cm] (ZI2) {$\zm$};
  \draw[->] (jI1) -- (ZI1);
  \draw[->] (jI2) -- (ZI2);
  \node[block, right of=ZI1] (DI1) {$\lift\bowtie$};
  \node[block, right of=ZI2] (DI2) {$\lift\bowtie$};
  \draw[->] (ZI1) -- (DI1);
  \draw[->] (ZI2) -- (DI2);
  \node[block, circle, below of=DI1, inner sep=0cm, node distance=.6cm] (sum) {$+$};
  \node[right of=sum] (output) {$o$};
  \draw[->] (ab) -- (sum);
  \draw[->] (DI1) -- (sum);
  \draw[->] (DI2) -- (sum);
  \draw[->] (a) -- (ab);
  \draw[->] (b) -- (ab);
  \draw[->] (a) -- (DI2);
  \draw[->] (b) -- (DI1);
  \draw[->] (sum) -- (output);
\end{tikzpicture}
\end{tabular}

Now we lift and incrementalize to get the circuit for $\inc{(\lift{\inc{(\lift{\bowtie})}})}$:

\begin{tikzpicture}[auto,>=latex]
  \node[] (a) {$a$}; 
  \node[below of=a, node distance=.8cm] (midway) {};
  \node[below of=ab, node distance=.8cm] (b) {$b$};

  \node[block, right of=a] (Ia) {$\I$};
  \node[block, right of=b] (Ib) {$\I$};
  \draw[->] (a) -- (Ia);
  \draw[->] (b) -- (Ib);
    
  \node[block, below of=Ia, node distance=.7cm] (ab) {$\lift\lift\bowtie$};
  \node[block, right of=Ia, node distance=1cm] (jI1) {$\lift\I$};
  \node[block, right of=Ib, node distance=1cm] (jI2) {$\lift\I$};
  \draw[->] (Ia) -- (jI1);
  \draw[->] (Ib) -- (jI2);
  \node[block, right of=jI1] (ZI1) {$\lift\zm$};
  \node[block, right of=jI2] (ZI2) {$\lift\zm$};
  \draw[->] (jI1) -- (ZI1);
  \draw[->] (jI2) -- (ZI2);
  \node[block, right of=ZI1] (DI1) {$\lift\lift\bowtie$};
  \node[block, right of=ZI2] (DI2) {$\lift\lift\bowtie$};
  \draw[->] (ZI1) -- (DI1);
  \draw[->] (ZI2) -- (DI2);
  \node[block, circle, below of=DI1, inner sep=0cm, node distance=.7cm] (sum) {$+$};
  \node[block, right of=sum] (D) {$\D$};
  \node[right of=D] (output) {$o$};
  \draw[->] (ab) -- (sum);
  \draw[->] (DI1) -- (sum);
  \draw[->] (DI2) -- (sum);
  \draw[->] (Ia) -- (ab);
  \draw[->] (Ib) -- (ab);
  \draw[->] (Ia) -- (DI2);
  \draw[->] (Ib) -- (DI1);
  \draw[->] (sum) -- (D);
  \draw[->] (D) -- (output);
\end{tikzpicture}

Applying the chain rule and the linearity of $\lift\I$ and $\lift\zm$ this becomes:

\begin{tikzpicture}[auto,>=latex, node distance=1.2cm]
  \node[] (a) {$a$}; 
  \node[block, below of=a, node distance=.8cm] (ab) {$\inc{(\lift\lift\bowtie)}$};
  \node[below of=ab, node distance=.8cm] (b) {$b$};
  \node[block, right of=a, node distance=1cm] (jI1) {$\lift{\I}$};
  \node[block, right of=b, node distance=1cm] (jI2) {$\lift{\I}$};
  \draw[->] (a) -- (jI1);
  \draw[->] (b) -- (jI2);
  \node[block, right of=jI1] (ZI1) {$\lift{\zm}$};
  \node[block, right of=jI2] (ZI2) {$\lift{\zm}$};
  \draw[->] (jI1) -- (ZI1);
  \draw[->] (jI2) -- (ZI2);
  \node[block, right of=ZI1] (DI1) {$\inc{(\lift\lift\bowtie)}$};
  \node[block, right of=ZI2] (DI2) {$\inc{(\lift\lift\bowtie)}$};
  \draw[->] (ZI1) -- (DI1);
  \draw[->] (ZI2) -- (DI2);
  \node[block, circle, below of=DI1, inner sep=0cm, node distance=.8cm] (sum) {$+$};
  \node[right of=sum] (output) {$o$};
  \draw[->] (ab) -- (sum);
  \draw[->] (DI1) -- (sum);
  \draw[->] (DI2) -- (sum);
  \draw[->] (a) -- (ab);
  \draw[->] (b) -- (ab);
  \draw[->] (a) -- (DI2);
  \draw[->] (b) -- (DI1);
  \draw[->] (sum) -- (output);
\end{tikzpicture}

We now have three applications of $\inc{(\lift\lift\bowtie)}$.  Each of these is the
incremental form of a bilinear operator, so it looks like in the end we will have $3\times3$ 
applications of $\lift\lift\bowtie$.  In fact, the overall expression can be simplified
(see~\cite{tr} for a precise derivation), and the end result only has 4 terms in $\lift\lift\bowtie$.

Here is the final form of the expanded join circuit:

\begin{tikzpicture}[auto,>=latex]
  \node[] (a) {$a$}; 
  \node[below of=a, node distance=.8cm] (b) {$b$};
  
  \node[block, right of=a] (LIa) {$\lift{\I}$};
  \node[block, above of=LIa, node distance=.8cm] (Ia) {$\I$};
  \node[block, right of=LIa] (IIa) {$\I$};
  \node[block, right of=Ia] (zIa) {$\zm$};
  \draw[->] (a) -- (LIa);
  \draw[->] (a) -- (Ia);
  \draw[->] (Ia) -- (zIa);
  \draw[->] (LIa) -- (IIa);
  
  \node[block, right of=b] (Ib) {$\I$};
  \node[block, below of=Ib, node distance=.8cm] (LIb) {$\lift\I$};
  \node[block, right of=Ib] (zb) {$\zm$};
  \node[block, right of=LIb] (IIb) {$\I$};
  \node[block, right of=IIb] (zIIb) {$\lift\zm$};
  \node[block, below of=IIb] (zIb) {$\lift\zm$};
  \draw[->] (b) -- (Ib);
  \draw[->] (b) -- (LIb);
  \draw[->] (Ib) -- (zb);
  \draw[->] (LIb) -- (IIb);
  \draw[->] (IIb) -- (zIIb);
  \draw[->] (LIb) -- (zIb);
  
  \node[block, right of=zIIb] (j1) {$\lift\lift\bowtie$};
  \node[block, above of=j1, node distance=.8cm]   (j2) {$\lift\lift\bowtie$};
  \node[block, above of=j2, node distance=.8cm]   (j3) {$\lift\lift\bowtie$};
  \node[block, above of=j3, node distance=.8cm]   (j4) {$\lift\lift\bowtie$};
  \draw[->] (zIIb) -- (j1);
  \draw[->] (a) -- (j1);
  \draw[->] (zb) -- (j2);
  \draw[->] (LIa) -- (j2);
  \draw[->] (IIa) -- (j3);
  \draw[->] (b) -- (j3);
  \draw[->] (zIa) -- (j4);
  \draw[->] (zIb) -- (j4);
  
  \node[block, right of=j3] (plus) {$+$};
  \draw[->] (j1) -- (plus);
  \draw[->] (j2) -- (plus);
  \draw[->] (j3) -- (plus);
  \draw[->] (j4) -- (plus); 
  \node[right of=plus] (o) {$o$};
  \draw[->] (plus) -- (o);
\end{tikzpicture}

Returning to $\inc{(\lift{\inc{(\lift\distinct)}})}$, we can compute its circuit by expanding
once using Proposition~\ref{prop-inc_distinct}:

\noindent
\begin{tabular}{m{3cm}m{.5cm}m{3.5cm}}
\begin{tikzpicture}[>=latex]
\node[] (input) {$i$};
\node[block, right of=input, node distance=1.5cm] (d) {$\inc{(\lift{\inc{(\lift{\distinct})}})}$};
\node[right of=d, node distance=1.5cm] (output) {$o$};
\draw[->] (input) -- (d);
\draw[->] (d) -- (output);
\end{tikzpicture}
& $\cong$ &
\begin{tikzpicture}[>=latex, node distance=.8cm]
    \node[] (input) {$i$};
    \node[block, right of=input] (I) {$\I$};
    \node[block, right of=I] (LI) {$\lift{\I}$};
    \node[block, right of=LI, node distance=1cm] (z) {$\lift{\zm}$};
    \node[block, below of=z, node distance=.8cm] (H) {$\lift{\lift{H}}$};
    \node[block, right of=H] (D) {$\D$};
    \node[right of=D] (output) {$o$};
    \draw[->] (input) -- (I);
    \draw[->] (I) -- node (mid) {} (LI);
    \draw[->] (LI) -- (z);
    \draw[->] (mid.center) |- (H);
    \draw[->] (z) -- (H);
    \draw[->] (H) -- (D);
    \draw[->] (D) -- (output);
\end{tikzpicture}
\end{tabular}

Finally, stitching all these pieces together we get the final circuit
shown in Figure~\ref{fig:recursive-example}.

\begin{figure*}[h]
\begin{tikzpicture}[>=latex]
  \node[] (Einput) {$\Delta$\code{E}};
  % generic part
  \node[block, right of=Einput, node distance=.8cm] (dE) {$\lift{\delta_0}$};
  \draw[->] (Einput) -- (dE); 
  \node[right of=dE] (empty) {};

  % join(a,b)
  \node[block, above of=empty, node distance=2.5cm] (b) { };
  \node[block, above of=b, node distance=.8cm] (a) { };
  \draw[->] (dE) -- (b);
  
  \node[block, right of=a] (LIa) {$\lift{\I}$};
  \node[block, above of=LIa, node distance=.8cm] (Ia) {$\I$};
  \node[block, right of=LIa] (IIa) {$\I$};
  \node[block, right of=Ia] (zIa) {$\zm$};
  \draw[->] (a) -- (LIa);
  \draw[->] (a) -- (Ia);
  \draw[->] (Ia) -- (zIa);
  \draw[->] (LIa) -- (IIa);
  
  \node[block, right of=b] (Ib) {$\I$};
  \node[block, below of=Ib, node distance=.8cm] (LIb) {$\lift\I$};
  \node[block, right of=Ib] (zb) {$\zm$};
  \node[block, right of=LIb] (IIb) {$\I$};
  \node[block, right of=IIb] (zIIb) {$\lift\zm$};
  \node[block, below of=IIb] (zIb) {$\lift\zm$};
  \draw[->] (b) -- (Ib);
  \draw[->] (b) -- (LIb);
  \draw[->] (Ib) -- (zb);
  \draw[->] (LIb) -- (IIb);
  \draw[->] (IIb) -- (zIIb);
  \draw[->] (LIb) -- (zIb);
  
  \node[block, right of=zIIb] (j1) {$\lift\lift\bowtie$};
  \node[block, above of=j1, node distance=.8cm]   (j2) {$\lift\lift\bowtie$};
  \node[block, above of=j2, node distance=.8cm]   (j3) {$\lift\lift\bowtie$};
  \node[block, above of=j3, node distance=.8cm]   (j4) {$\lift\lift\bowtie$};
  \draw[->] (zIIb) -- (j1);
  \draw[->] (a) -- (j1);
  \draw[->] (zb) -- (j2);
  \draw[->] (LIa) -- (j2);
  \draw[->] (IIa) -- (j3);
  \draw[->] (b) -- (j3);
  \draw[->] (zIa) -- (j4);
  \draw[->] (zIb) -- (j4);
  
  \node[block, right of=j3, circle, inner sep=0cm] (plus) {$+$};
  \draw[->] (j1) -- (plus);
  \draw[->] (j2) -- (plus);
  \draw[->] (j3) -- (plus);
  \draw[->] (j4) -- (plus); 
  
  % relational query
  \node[block, right of=plus] (pj) {$\lift{\lift{\pi}}$};
  \node[block, below of=empty] (p1) {$\lift{\lift{\pi}}$};
  \node[block, right of=p1] (s1) {$\lift{\lift{\sigma}}$};
  \node[block, below of=p1, node distance=.6cm] (p2)  {$\lift{\lift{\pi}}$};
  \node[block, right of=p2] (s2) {$\lift{\lift{\sigma}}$};
  \node[below of=pj, node distance=.6cm] (mid) {};
  \node[block, circle, right of=empty, inner sep=0cm, node distance=8cm] (relplus) {$+$};
  
  \draw[->] (plus) -- (pj);
  \draw[->] (dE) -- (p1);
  \draw[->] (p1) -- (s1);
  \draw[->] (dE) -- (p2);
  \draw[->] (p2) -- (s2);
  \draw[->] (dE) -- (relplus);
  \draw[->] (pj) -- (relplus);
  \draw[->] (s1) -- (relplus);
  \draw[->] (s2) -- (relplus);
  
    % distinct
    \node[block, right of=relplus] (distI) {$\I$};
    \node[block, right of=distI] (distLI) {$\lift{\I}$};
    \node[block, right of=distLI, node distance=1cm] (distz) {$\lift{\zm}$};
    \node[block, below of=distz, node distance=.8cm] (distH) {$\lift{\lift{H}}$};
    \node[block, right of=distH] (distD) {$\D$};
    \draw[->] (relplus) -- (distI);
    \draw[->] (distI) -- node (distmid) {} (distLI);
    \draw[->] (distLI) -- (distz);
    \draw[->] (distmid.center) |- (distH);
    \draw[->] (distz) -- (distH);
    \draw[->] (distH) -- (distD);
  
  % generic part
  \node[block, right of=distD] (S) {$\lift{\int}$};
  \node[right of=S, node distance=.8cm] (output)  {$\Delta$\code{R}};
  \draw[->] (distD) -- node (o) {} (S);
  \draw[->] (S) -- (output);
  \node[block, above of=a, node distance=1.2cm] (z) {$\lift{\zm}$};
  \draw[->] (o.center) |- (z);
  \draw[->] (z) -- (a);
\end{tikzpicture}
\caption{Final form of circuit from \refsec{sec:recursive-example}\label{fig:recursive-example}.}
\end{figure*}

\subsection{Complexity of incremental recursive queries}

\paragraph{Time complexity}

The time complexity of an incremental recursive query can be estimated as a product of
the number of fixed point iterations and the complexity of each iteration. The
incrementalized circuit (\ref{eq:increcursive}) performs the same number of
iterations as the non-incremental circuit (\ref{eq:seminaive}) in the worst case:
once the non-incremental circuit reaches the fixed point, its output is constant
and so is its derivative computed by the incrementalized circuit.

Consider a nested stream of changes $s \in \stream{\stream{A}}, s[t_1][t_2]$,
where $t_1$ is the input timestamp and $t_2$ is the fixed point iteration number.
The unoptimized loop body $\inc{(\lift{\inc{(\lift{T})}})} = 
\D \circ \lift{\D} \circ \lift{\lift{T}}
\circ \lift{\I} \circ \I$ has the same time complexity as $T$ applied to the
aggregated input of size $R(s)[t_1][t_2] \defn \norm{(\lift{\I} \circ
\I)(s)[t_1][t_2]} = \norm{\sum_{(i_1,i_2) \leq (t_1, t_2)} s[i_1][i_2]}$.  As
before, an optimized circuit can be significantly more efficient.  For instance,
by applying Theorem~\ref{bilinear} twice, to $\bowtie$ and $\lift{\bowtie}$, we
obtain a circuit for nested incremental join $s_1
\inc{(\lift{\inc{(\lift{\bowtie})}})} s_2$ that runs in
$O(\norm{\lift{\I}(s1)[t1][t2]} \times \norm{\I(s2)[t1][t2]}) \ll 
O(R(s_1) \times R(s_2))$ (because each term is correspondingly smaller).

\paragraph{Space complexity} Integration ($\I$) and differentiation ($\D$) of a
stream $s \in \stream{\stream{A}}$ uses memory proportional to
$\sum_{t_2}\norm{\sum_{t_1}s[t_1][t_2]}$, i.e., the total size of changes
aggregated over columns of the matrix.  The unoptimized circuit integrates
and differentiates respectively inputs and outputs of the recursive program
fragment.  As we move $\I$ and $\D$ inside the circuit using the chain rule, we
additionally store changes to intermediate streams.  Effectively we cache results of 
fixed point iterations from earlier timestamps to update them efficiently as new input changes arrive.
Notice that space is proportional to the number of iterations of the inner while loop.


